\documentclass[11pt,a4paper]{article}

\usepackage[margin=1in]{geometry}
\usepackage{hyperref}
\usepackage{enumitem}
\usepackage{xcolor}
\usepackage{titling}

\makeatletter
\newcommand{\BvrSetLabInstructor}[1]{%
  \gdef\Bvr@LabInstructor{#1}}
\newcommand{\Bvr@LabInstructor}{Dr. Paramveer Kaur}
\newcommand{\BvrLabInstructorValue}{\Bvr@LabInstructor}
\makeatother

\pretitle{\begin{center}\bfseries\fontsize{18bp}{18bp}\selectfont}

\makeatletter
\newcommand{\BvrSubtitle}[1]{%
  \posttitle{%
    \par\end{center}
    \begin{center}\large#1\end{center}
    \vskip 0.5em}%
}
\makeatother

\newcommand{\BvrHint}[1]{%
  {\normalfont\normalsize\itshape\hfill #1}}

\title{Uber Lite Ride Sharing System}

\BvrSetLabInstructor{Dr. Paramveer Kaur}

\BvrSubtitle{%
  Project Proposal for UCS503P  \\
  Submitted to: \BvrLabInstructorValue \\ \bfseries
  Thapar Institute of Engineering and Technology%
}

\author{
  Ayush, CSED \\
  \texttt{1024030740, aayush\_be24@thapar.edu}
  \\[0.5em]
  Gurleen Kaur, CSED \\
  \texttt{1024030325, gkaur6\_be24@thapar.edu}
  \\[0.5em]
  Anshika, CSED \\
  \texttt{1024030749, aanshika\_be24@thapar.edu}
}

\date{\today}

\begin{document}

\maketitle


\section{Higher-order goal%
  \BvrHint{\ldots secondary framing}}

This project aims to improve the accessibility and efficiency of
urban ride-sharing by reducing the time, cost, and effort required
for users to find and book a suitable ride. The broader goal is to
provide a lightweight and practical transportation system that
connects passengers with available drivers while maintaining a
simple and responsive user experience.

The immediate engineering objective is to translate the core
ride-sharing workflow---from requesting a ride to matching a driver
and completing the trip---into a measurable and deployable
full-stack web application. The system will focus on practical
aspects such as location-based driver matching, ride booking,
fare estimation, trip-status management, authentication, and
reliable communication through REST APIs.


\section{Time-to-value%
  \BvrHint{\ldots primary heuristic}}

\begin{itemize}[leftmargin=0.7cm]

    \item We will deliver meaningful functionality quickly by
    prototyping the core ride-booking workflow first: user
    registration, authentication, ride request, driver matching,
    booking confirmation, and ride-status updates.

    \item We will prioritize fast feedback over unnecessary
    complexity by implementing the system in short iterations,
    validating the main passenger and driver workflows early, and
    improving the system based on testing results.

    \item The first iteration will focus on a working end-to-end
    ride-sharing workflow rather than advanced features. This allows
    the project to demonstrate measurable value early while leaving
    room for later improvements such as optimized driver matching,
    improved database queries, fare adjustments, real-time ride
    updates, feedback and rating functionality, and support for
    multiple concurrent ride requests.

    \item Success will be defined by what can be delivered and
    measured in the first iteration: a user should be able to
    authenticate, request a ride, the system should identify an
    appropriate available driver, and the ride should progress
    through clearly defined states.

\end{itemize}


\section{Problem Statement}

Traditional transportation services can be inconvenient and difficult
to manage when users have to search for available vehicles, contact
drivers manually, and coordinate ride details through separate
channels. Users may also face difficulties in knowing whether a
driver is available, tracking the status of a ride request, and
estimating the expected fare before or during a trip.

The main problems addressed by UberLite include:

\begin{itemize}

    \item \textbf{Difficulty in finding available drivers:} Users may
    have to manually search for or contact drivers, leading to delays
    in booking a ride.

    \item \textbf{Lack of centralized ride management:} Ride
    requests, driver availability, and ride details are difficult to
    manage without a unified platform.

    \item \textbf{Limited ride status visibility:} Users may not have
    clear information about whether a ride is requested, accepted,
    in progress, or completed.

    \item \textbf{Uncertainty in ride information:} Users may not have
    a simple way to view important ride details such as pickup
    location, destination, driver information, and estimated fare.

    \item \textbf{Manual management of ride requests:} Handling
    bookings and updating ride information manually can result in
    delays, inconsistencies, and errors.

\end{itemize}

These issues create the need for a simple and centralized system that
can manage the basic interaction between passengers and drivers while
providing clear ride information and status updates. UberLite
addresses this need through a lightweight digital platform for
requesting, matching, and managing rides.


\section{Proposed Solution}

\textit{...Software Proposition}

\subsection{Overview}

Build a lightweight full-stack ride-sharing web application using
React, Node.js, Express.js, PostgreSQL, and Prisma. The system will
provide a React-based frontend for passengers and drivers and a
Node.js and Express.js backend for application logic and REST API
communication. The system will provide the following core
functionality:

\begin{itemize}

    \item \textbf{User registration and authentication:} Passengers
    and drivers can register and authenticate before accessing
    role-specific functionality.

    \item \textbf{Passenger and driver dashboards:} The React
    frontend will provide role-specific dashboards for passengers
    and drivers, allowing each user type to access the functions
    relevant to their workflow.

    \item \textbf{Ride request module:} Passengers can provide pickup
    and destination locations and submit a ride request through the
    web application.

    \item \textbf{Driver availability and matching:} The system
    maintains driver availability and matches ride requests with
    suitable available drivers based on location.

    \item \textbf{Ride management workflow:} Drivers can accept or
    reject ride requests and update the ride status throughout the
    journey.

    \item \textbf{Fare estimation and calculation:} The system
    calculates an estimated and final fare using predefined pricing
    rules based on ride-related information.

    \item \textbf{Ride status tracking:} The system maintains the
    current status of each ride, such as Requested, Accepted,
    In Progress, or Completed.

    \item \textbf{Ride history:} Completed rides and their relevant
    details are stored so that previous trips can be retrieved.

\end{itemize}

The frontend will communicate with the Express.js backend through
REST APIs and structured JSON data. PostgreSQL will be used for
persistent storage of users, drivers, rides, locations, fare details,
and ride history, with Prisma providing the database access layer.


\subsection{Core Workflow}

\begin{enumerate}

    \item Passenger registers and authenticates with the system.

    \item Passenger submits a ride request containing pickup and
    destination information.

    \item System validates the request and checks for available
    drivers.

    \item System identifies a suitable driver using the driver
    matching mechanism.

    \item Driver accepts or rejects the ride request.

    \item Once accepted, the ride status changes from
    \textit{Requested} to \textit{Accepted}.

    \item Driver starts the trip and the ride status changes to
    \textit{In Progress}.

    \item Driver completes the trip and the system calculates the
    final fare.

    \item The completed ride is stored in the ride history.

    \item After the ride is completed, the passenger can provide
    a rating and optional feedback for the driver.

\end{enumerate}


\subsection{Operational Constraints for an Educational Setting}

\begin{itemize}

    \item Keep the application lightweight and suitable for
    deployment on commonly available computing resources.

    \item Focus on the essential ride-booking workflow rather than
    replicating the complete functionality of commercial
    ride-hailing platforms.

    \item Use a modular full-stack architecture so that frontend,
    backend, and database components can be tested and extended
    independently.

    \item Use mature technologies and libraries rather than
    developing infrastructure from scratch.

    \item Avoid unnecessary dependence on complex external services
    during the initial implementation.

    \item Prioritize correctness, reliability, efficient data
    management, and scalability of the core application.

\end{itemize}


\section{Solution Approach}
\BvrHint{\ldots Engineering focus}

This is an engineering course project, so the focus will be on
building a reliable, modular, and scalable full-stack application
rather than developing novel algorithms. The implementation will
prioritize the core passenger-driver workflow, efficient driver
matching, well-defined ride-status transitions, secure API access,
and incremental development.


\subsection{Backend Architecture}
\BvrHint{\ldots Node.js and Express.js REST-based solution}

The system will use a modular Node.js and Express.js backend exposed
through REST APIs. The major components will be separated according
to their responsibilities:

\begin{itemize}

    \item \textbf{Authentication Module:} Handles user registration,
    login, JWT-based authentication, and role-based access for
    passengers and drivers.

    \item \textbf{Ride Management Module:} Creates ride requests,
    manages ride information, and controls valid ride-status
    transitions.

    \item \textbf{Driver Management Module:} Maintains driver
    availability, location information, and active ride assignments.

    \item \textbf{Matching Module:} Searches for suitable available
    drivers based on their location relative to the pickup point.

    \item \textbf{Fare Module:} Calculates estimated and final fares
    using predefined pricing rules.

    \item \textbf{Feedback Module:} Records passenger ratings and
    optional feedback for completed rides and associates them with
    the relevant driver and ride.

    \item \textbf{Data Management Module:} Handles persistent storage
    and retrieval of users, drivers, rides, feedback, and ride history
    using PostgreSQL and Prisma.

    \item \textbf{Validation Module:} Validates incoming API
    requests before processing them to maintain consistent and
    reliable application data.

\end{itemize}

The Express.js backend will expose these functionalities through
REST endpoints, with JSON used for structured request and response
data.


\subsection{Driver Matching}
\BvrHint{\ldots lightweight matching}

Driver matching will initially use a simple location-based approach:

\begin{itemize}

    \item Maintain the availability and current location of
    registered drivers.

    \item When a passenger requests a ride, identify currently
    available drivers.

    \item Calculate or estimate the distance between the pickup
    location and available drivers.

    \item Prioritize the nearest suitable available driver.

    \item Update the driver's availability after assignment so that
    the same driver cannot be assigned to multiple active rides.

\end{itemize}

The initial implementation will use a straightforward distance-based
matching strategy. The matching module will remain independent from
the rest of the system so that it can later be optimized using more
efficient search or location-based techniques.


\subsection{Route and Distance Calculation}
\BvrHint{\ldots practical location-based computation}

The system will use location and distance calculations where
required for estimating the distance between relevant locations.
The calculated distance can be used by the fare calculation module
and can also support driver-selection decisions.

The distance calculation component will remain separate from the
ride-management logic so that it can be improved independently as
the system grows.


\subsection{Ride Workflow Management}
\BvrHint{\ldots controlled state transitions}

The system will manage every ride through a clearly defined workflow:

\begin{center}
\textit{Requested} $\rightarrow$
\textit{Accepted} $\rightarrow$
\textit{In Progress} $\rightarrow$
\textit{Completed}
\end{center}

Only valid state transitions will be permitted. This ensures that
ride information remains consistent and prevents invalid operations,
such as completing a ride that has not been accepted or started.


\subsection{Data Management}
\BvrHint{\ldots structured and persistent data}

The system will maintain structured records for passengers, drivers,
rides, and feedback. PostgreSQL will be used as the relational
database for persistent storage of users, drivers, rides, locations,
fare details, feedback, ratings, and ride history.

Prisma will be used as the database access layer and ORM for defining
the application data models and interacting with PostgreSQL. The
database layer will be separated from the core ride-management logic
so that storage-related changes do not require redesigning the
complete application.

Appropriate indexes and relational constraints will be considered
for frequently queried and related data such as driver availability,
ride status, user roles, and ride history.


\subsection{Authentication and Security}
\BvrHint{\ldots secure API access}

Authentication will be implemented so that passengers and drivers
can access only the operations relevant to their roles.

JWT-based token authentication will be used for API access. Input
validation will be performed before processing requests to prevent
invalid or inconsistent data from entering the system.

Zod will be used for request validation, while CORS will be
configured to allow controlled communication between the React
frontend and the Express.js backend.


\subsection{Fare Calculation}
\BvrHint{\ldots deterministic pricing}

The fare calculation will use a predefined pricing model rather than
dynamic pricing. A simple model can be represented as:

\[
\text{Fare} =
\text{Base Fare} +
(\text{Distance} \times \text{Rate per km})
\]

The same calculation rules will be applied consistently so that fare
estimates and final fares are predictable and easy to test.


\section{Evaluation Criteria}
\BvrHint{\ldots measurable and attributable}

The success of UberLite will be evaluated using measurable metrics
directly related to the core ride-booking workflow. The evaluation
will focus on system correctness, efficiency, reliability,
application performance, and successful completion of important
frontend workflows rather than comparison with commercial
ride-hailing platforms.


\subsection{Primary Evaluation Metric}
\BvrHint{\ldots ride request to driver assignment}

\textbf{Driver Assignment Time (DAT):} The time taken between a
passenger submitting a valid ride request and the system assigning
an available driver.

\begin{itemize}

    \item \textbf{Target:} The median driver assignment time should
    remain below a predefined threshold during testing.

    \item \textbf{Measurement:} The metric will be calculated using
    timestamps recorded when the ride request is created and when a
    driver is successfully assigned.

    \item \textbf{Attribution:} Both timestamps will be generated and
    stored by the backend, allowing the metric to be measured directly
    from application logs and database records.

\end{itemize}


\subsection{Secondary Evaluation Metrics}

\begin{itemize}

    \item \textbf{Ride Request Success Rate:} Percentage of valid
    ride requests that are successfully matched with an available
    driver.

    \item \textbf{Matching Accuracy:} Percentage of ride requests
    for which the system selects the intended nearest or most
    suitable available driver according to the defined matching rule.

    \item \textbf{Fare Calculation Accuracy:} Percentage of test
    rides for which the calculated fare matches the expected fare
    according to the predefined pricing formula.

    \item \textbf{Workflow Correctness:} Percentage of tested rides
    that successfully follow the valid sequence of states from
    \textit{Requested} to \textit{Accepted}, \textit{In Progress},
    and finally \textit{Completed}.

    \item \textbf{API Response Time:} Average time taken by important
    backend operations such as authentication, ride creation,
    driver matching, and ride-status updates.

    \item \textbf{Frontend Workflow Completion:} Percentage of test
    scenarios in which passengers and drivers can successfully
    complete their intended workflows through the React interface.

    \item \textbf{Feedback Submission Reliability:} Percentage of
    eligible completed rides for which ratings and feedback can be
    successfully submitted and stored.

    \item \textbf{Reliability:} Percentage of test scenarios
    completed without API failures, invalid state transitions,
    inconsistent ride assignments, or system errors.

\end{itemize}


\subsection{Pilot Validation Plan}
\BvrHint{\ldots high-level}

The system will be validated using controlled test scenarios that
represent the main passenger and driver workflows.

\begin{itemize}

    \item Test the system with multiple passenger and driver accounts
    representing realistic ride-booking scenarios.

    \item Measure driver assignment time for a set of ride requests
    under different numbers of available drivers.

    \item Test cases where drivers accept requests, reject requests,
    or become unavailable before assignment.

    \item Verify that invalid ride-status transitions are rejected by
    the system.

    \item Compare calculated fares against manually computed expected
    values using the predefined fare formula.

    \item Test API endpoints for valid requests, invalid input, and
    unauthorized access.

    \item Test the passenger and driver dashboards for successful
    navigation, authentication, ride-request submission, and
    ride-status display.

    \item Test rating and feedback submission after completed rides
    and verify that the feedback is correctly associated with the
    relevant ride and driver.

\end{itemize}

The collected results will be used to identify performance issues,
workflow errors, and reliability problems. Improvements will then be
made iteratively based on the observed results.


\section{Scalability}
\BvrHint{\ldots with foresight}

Although UberLite is initially developed as an academic full-stack
application, the system will be designed to handle growth in users,
drivers, and ride requests without requiring a complete redesign.

\subsection{Efficient Driver Matching}
\BvrHint{\ldots reducing unnecessary search}

Driver matching is the main scalability concern. Instead of
repeatedly checking every driver for every ride request, the matching
logic will remain modular so that more efficient location-based
search techniques can be introduced if the number of drivers
increases significantly.

The initial implementation will use a straightforward distance-based
matching strategy suitable for a controlled academic environment.


\subsection{Data and Backend Scalability}
\BvrHint{\ldots efficient operations}

User, driver, ride, and feedback information will be stored in
PostgreSQL with appropriate indexing and relational constraints for
frequently queried and related data. The Express.js backend will be
structured into separate routes, controllers, validation, business
logic, and database operations.

Prisma will provide a consistent database access layer between the
backend application and PostgreSQL. This separation will allow
individual components to be optimized without requiring major
changes to the complete application.


\subsection{Scalability Path}
\BvrHint{\ldots prototype to larger system}

The initial system will target a controlled academic environment.
For larger deployments, scalability can be improved through more
efficient driver-search strategies, optimized PostgreSQL queries
and indexes, better handling of concurrent ride requests,
pagination for large datasets, and multiple backend instances when
required.


\section{Technology and Engine Availability}
\BvrHint{\ldots mature engineering components}

UberLite will use a mature set of technologies and libraries that
are well suited for implementing a full-stack web application and
REST API system. These components allow the project to focus on
ride-booking workflow, correctness, scalability, and evaluation
rather than developing basic infrastructure from scratch.

\begin{itemize}

    \item \textbf{React:} Used to develop the frontend interface for
    passengers and drivers and provide role-specific dashboards for
    the ride-booking workflow.

    \item \textbf{Node.js:} Used as the server-side JavaScript
    runtime for implementing the backend application.

    \item \textbf{Express.js:} Used for implementing REST APIs and
    handling HTTP requests, responses, and backend application logic.

    \item \textbf{PostgreSQL:} Used as the relational database for
    persistent storage of users, drivers, rides, locations, fare
    details, feedback, ratings, and ride history.

    \item \textbf{Prisma:} Used as the ORM and database access layer
    for defining application data models and interacting with
    PostgreSQL from the Node.js backend.

    \item \textbf{Zod:} Used for request-body validation and ensuring
    that API inputs follow the required structure.

    \item \textbf{JWT:} Used to implement token-based authentication
    and protect role-specific API endpoints.

    \item \textbf{CORS:} Used to control communication between the
    React frontend and the Express.js backend.

    \item \textbf{Postman:} Used to test and validate REST API
    endpoints during development.

    \item \textbf{Git/GitHub:} Used for version control,
    collaborative development, and maintaining the project
    repository.

\end{itemize}

These technologies provide a mature development environment for the
project and allow the team to focus on the engineering aspects of
UberLite, particularly efficient driver matching, ride workflow
management, correctness, testing, and scalability.


\section{Project Scope and Deliverables}

The project will be developed incrementally, with the first iteration
focused on implementing and demonstrating the complete core
ride-booking workflow through the backend APIs and web-based
frontend.

\subsection{Initial Deliverable}
\BvrHint{\ldots first iteration}

\begin{itemize}

    \item User registration and authentication for passengers and
    drivers.

    \item React-based passenger and driver dashboards.

    \item Ride request creation with pickup and destination
    information.

    \item Driver availability management and basic driver matching.

    \item Ride workflow from \textit{Requested} to
    \textit{Accepted}, \textit{In Progress}, and
    \textit{Completed}.

    \item Basic fare estimation and final fare calculation.

    \item Location and distance calculation for driver matching and
    fare estimation.

    \item Persistent storage of users, drivers, rides, and ride
    history using PostgreSQL.

    \item REST API endpoints for the core system functionality.

    \item Input validation and role-based authentication.

    \item Unit and integration testing of the main backend modules.

\end{itemize}


\subsection{Subsequent Deliverables}
\BvrHint{\ldots iteration-friendly}

\begin{itemize}

    \item Improved driver-matching and search efficiency.

    \item Enhanced handling of concurrent ride requests.

    \item Improved ride-history search and filtering.

    \item Rating and feedback functionality for completed rides.

    \item Real-time ride-status updates between passenger and driver
    interfaces.

    \item Periodic driver location updates and improved ride tracking.

    \item Additional validation, authentication, and security
    improvements.

    \item Improved frontend usability and responsive design.

    \item Performance testing and scalability evaluation.

    \item Deployment of the application for online demonstration,
    subject to project time and available resources.

\end{itemize}


\section{Risks and Mitigations}

\begin{itemize}

    \item \textbf{Driver matching efficiency:} A simple search may
    become inefficient as the number of drivers increases. This will
    be addressed through a modular matching component that can be
    optimized later.

    \item \textbf{Concurrent ride requests:} Multiple requests may
    attempt to assign the same driver. Driver availability will be
    verified before assignment and updated immediately after a
    successful assignment.

    \item \textbf{Data consistency:} Invalid ride states or
    inconsistent records may occur if operations are not validated.
    The system will enforce defined ride-state transitions and
    validate API inputs. PostgreSQL transactions and relational
    constraints can also be used where required.

    \item \textbf{Security:} Unauthorized users may attempt to access
    restricted operations. JWT-based authentication and role-based
    access control will be applied to relevant API endpoints.

    \item \textbf{Frontend-backend communication:} Incorrect API
    requests or cross-origin configuration may cause communication
    failures. REST API contracts, request validation, and appropriate
    CORS configuration will be used to reduce such issues.

    \item \textbf{Feedback feature complexity:} Rating and feedback
    functionality will be implemented only for completed rides and
    kept lightweight so that it does not interfere with the core
    ride-booking workflow.

    \item \textbf{Real-time feature complexity:} Real-time ride
    updates and driver-location tracking may require additional
    implementation effort. These features will be treated as
    subsequent deliverables and prioritized based on available
    development time after the core workflow is stable.

    \item \textbf{Evaluation ambiguity:} Performance and correctness
    can be difficult to assess without clearly defined metrics.
    Backend timestamps, test cases, and measurable evaluation
    criteria will be established before final testing.

    \item \textbf{Project complexity:} Adding too many commercial
    ride-hailing features may make the project difficult to complete.
    The initial implementation will remain focused on the core
    ride-booking workflow, with additional features added only when
    the core system is stable.

\end{itemize}


\section{Summary}

UberLite proposes a lightweight full-stack ride-sharing application
using React, Node.js, Express.js, PostgreSQL, and Prisma. The system
provides the essential workflow of a ride-hailing system, including
user authentication, passenger and driver dashboards, ride requests,
driver matching, location and distance calculation, fare calculation,
ride-status management, ride history, and post-ride feedback.

The project is designed as an engineering-focused solution rather
than a research-driven system. It emphasizes time-to-value through
incremental development, measurable evaluation through driver
assignment time and workflow correctness, and scalability through
efficient driver matching, structured relational data management,
and modular frontend and backend design.

The initial implementation will provide a complete and testable
ride-booking workflow through a React frontend and Node.js/Express.js
backend, with PostgreSQL providing persistent data storage and
Prisma providing the database access layer. Subsequent development
may extend the system with rating and feedback functionality,
real-time ride-status updates, driver location tracking, improved
frontend interaction, and online deployment depending on the
available project time and testing results.


\end{document}
